\uuid{F2tG}
\titre{DL en l'infini}
\auteur{hanani ; mzari}
\organisation{exo7}
\licence{CC BY-NC-SA 4.0}
\datecreate{2020-01-20}
\niveau{L1}
\module{Analyse}
\chapitre{Développement limité}
\sousChapitre{Calculs}
\difficulte{3}
\theme{Opérations sur les DL, DL en l'infini}
\exo7qcmid{lille-2.developpements-limites.014}

\contenu{
\texte{
Soit $a\in \Rr$ et $f$ une fonction. On dit que $f(x)$ est un $o\left((x-a)^n\right)$ au voisinage de $a$ et on écrit $f(x)=o((x-a)^n)$ si $\displaystyle \lim _{x\to a}\frac{f(x)}{(x-a)^n}=0$. On note par $DL_n(a)f(x)$ le développement limité de $f$ à l'ordre $n$ en $a$. On note aussi $DL_n(a^+)f(x)$ (resp. $DL_n(a^-)f(x)$) le développement limité de $f$ à l'ordre $n$ à droite en $a$ (resp. à gauche en $a$).
}

\begin{question}
\qclassification{125.02}
\qtags{motcle=DL en l'infini}
Les égalités suivantes portent sur des développements limités au voisinage de $+\infty$. Cocher celles qui sont vraies :
\begin{answers}
\good{$\displaystyle DL_3(+\infty)\sin \frac{1}{x}=\frac{1}{x}-\frac{1}{6x^3}+o\left(\frac{1}{x^3}\right)$.}
\bad{$\displaystyle DL_2(+\infty)\frac{1}{1+x^2}=1-\frac{1}{x^2}+o\left(\frac{1}{x^2}\right)$.}
\good{$\displaystyle DL_2(+\infty)\ln \frac{x+1}{x}=\frac{1}{x}-\frac{1}{2x^2}+o\left(\frac{1}{x^2}\right)$.}
\bad{$\displaystyle DL_2(+\infty)\frac{1}{1+x}=1-\frac{1}{x}+\frac{1}{x^2}+o\left(\frac{1}{x^2}\right)$.}
\end{answers}
\vskip2mm
\begin{explanations}
Pour écrire un $DL(+\infty)$ de $f(x)$, on écrit un $DL(0^+)$ de $g(t)=f(1/t)$. Ainsi
$$\begin{array}{l}\displaystyle \sin\left(t\right)=t-\frac{t^3}{6}+o(t^3)\Rightarrow DL_3(+\infty)\sin \frac{1}{x}=\frac{1}{x}-\frac{1}{6x^3}+o\left(\frac{1}{x^3}\right).\\ \\ \displaystyle \frac{1}{1+1/t^2}=\frac{t^2}{1+t^2}=t^2+o(t^2)\Rightarrow DL_2(+\infty)\frac{1}{1+x^2}=\frac{1}{x^2}+o\left(\frac{1}{x^2}\right).\\ \\ \displaystyle \ln \left(\frac{1/t+1}{1/t}\right)=\ln (1+t)=t-\frac{t^2}{2}+o(t^2)\Rightarrow DL_2(+\infty)\ln \frac{x+1}{x}=\frac{1}{x}-\frac{1}{2x^2}+o\left(\frac{1}{x^2}\right).\\ \\ \displaystyle \frac{1}{1+1/t}=\frac{t}{1+t}=t-t^2+o(t^2)\Rightarrow DL_2(+\infty)\frac{1}{1+x}=\frac{1}{x}-\frac{1}{x^2}+o\left(\frac{1}{x^2}\right).\end{array}$$
\end{explanations}
\end{question}
}
